% ATS-friendly resume: single column, no tables, no icons, no graphics,
% standard section headings, plain-text contacts, selectable text.
\documentclass[10pt,a4paper]{article}

\usepackage[T2A]{fontenc}
\usepackage[utf8]{inputenc}
\usepackage[russian]{babel}
\usepackage[left=1.5cm,right=1.5cm,top=1.3cm,bottom=1.3cm]{geometry}
\usepackage{enumitem}
\usepackage[hidelinks]{hyperref}

\pagestyle{empty}
\hyphenpenalty=10000
\exhyphenpenalty=10000
\raggedright
\setlength{\parindent}{0pt}
\setlength{\parskip}{0pt}

\newcommand{\cvsection}[1]{%
  \vspace{7pt}{\large\bfseries #1}\par\vspace{2pt}\hrule\vspace{4pt}}

% #1 -- должность, #2 -- даты, #3 -- организация и город
\newcommand{\cventry}[3]{%
  \textbf{#1} \textbar{} #3 \textbar{} #2\par\vspace{2pt}}

\newlist{cvitems}{itemize}{1}
\setlist[cvitems]{label=\textbullet, leftmargin=1.2em, topsep=2pt, itemsep=1pt, parsep=0pt}

\begin{document}

{\LARGE\bfseries Кирилл Козлов}\par\vspace{3pt}
ML-разработчик / Python-разработчик\par\vspace{3pt}
Москва, Россия | Телефон: +7 (910) 968 1153 | Email: kirillandreevickozlov@yandex.ru\par
GitHub: github.com/noobmaster-ru | Telegram: @noobmaster\_rus

\cvsection{Профиль}
Python- и ML-разработчик, выпускник ВМК МГУ. Разрабатываю LLM-агентов и RAG-системы
(Retrieval-Augmented Generation) в X5 Tech: бэкенд на FastAPI, интеграция больших языковых
моделей (LLM) и распознавания документов (OCR), развёртывание в Kubernetes через GitLab CI/CD.
Экономический эффект внедрённых решений --- 8 и 10 FTE.

\cvsection{Ключевые навыки}
\textbf{Языки программирования:} Python\par
\textbf{Искусственный интеллект и ML:} большие языковые модели (LLM), LLM-агенты,
RAG (Retrieval-Augmented Generation), OpenAI API, LangChain, промпт-инжиниринг,
распознавание документов (OCR), Claude Code\par
\textbf{Фреймворки:} FastAPI, SQLAlchemy (async), Alembic, Pydantic, pytest\par
\textbf{Библиотеки:} asyncio, httpx, aiohttp, pandas, NumPy, scikit-learn, structlog\par
\textbf{Базы данных:} PostgreSQL, Redis\par
\textbf{Инфраструктура и DevOps:} Docker, Kubernetes, Helm, GitLab CI/CD, Apache Kafka, Git, Linux\par
\textbf{Языки:} русский --- родной, английский --- Intermediate (B1)

\cvsection{Опыт работы}
\cventry{ML-разработчик}{Апрель 2026 --- настоящее время}{X5 Tech, департамент ИИ-агентов, Москва}
\begin{cvitems}
    \item Разработал LLM-агента для анализа юридических документов; экономический эффект для X5 Group --- 8 FTE.
    \item Развернул сервис в кластере Kubernetes в средах dev и prod с использованием Helm и GitLab CI/CD.
    \item Разработал бэкенд на FastAPI для сервиса распознавания коммунальных платежей (OCR) по помещениям X5 Group в составе команды ML-инженеров; экономический эффект --- 10 FTE.
    \item Автоматизирую анализ иностранных инвойсов и счетов бизнес-единицы X5 Import с помощью OCR и больших языковых моделей.
\end{cvitems}

\cvsection{Проекты}
\cventry{Python-разработчик, AxiomAI}{Октябрь 2025 --- Апрель 2026}{github.com/noobmaster-ru/axiomai, Москва}
\begin{cvitems}
    \item Разработал и развернул Telegram-бота на aiogram с использованием OpenAI API, PostgreSQL, Redis и Docker.
    \item Снизил нагрузку на менеджеров на 80\% и повысил скорость ответов в бизнес-коммуникациях.
\end{cvitems}
\vspace{4pt}

\cventry{Python-разработчик, WB Parser (Telegram Mini App)}{Июль 2025 --- Сентябрь 2025}{Москва}
\begin{cvitems}
    \item Разработал парсер поисковой выдачи Wildberries с фильтрацией по органической позиции товара.
    \item Упаковал продукт в Telegram Mini App с бэкендом и базой данных; привлёк 220 активных пользователей.
\end{cvitems}
\vspace{4pt}

\cventry{Python-разработчик, Placa}{Июль 2024 --- Декабрь 2024}{Москва}
\begin{cvitems}
    \item Разработал сервис аналитики в Google Sheets для продавцов Wildberries на Google Apps Script с интеграцией Wildberries API.
    \item Автоматизировал ежедневную загрузку данных: копирование шаблона таблицы, ввод API-токена, запуск с панели управления.
    \item Обеспечил около 300 продавцов ежедневными данными о рекламных расходах.
\end{cvitems}

\cvsection{Образование}
\cventry{МГУ имени М.В. Ломоносова}{Август 2022 --- Август 2026}{Факультет вычислительной математики и кибернетики (ВМК), Москва}
\begin{cvitems}
    \item Deep Learning School (DLS), направления Computer Vision и NLP --- Школа глубокого обучения МФТИ.
    \item Летняя школа по A/B-тестированию, X5 Tech.
\end{cvitems}

\cvsection{Интересы}
Волейбол, искусственный интеллект, тренажёрный зал, экономика, история.

\end{document}
